\documentclass[12pt]{article}

\usepackage[margin=1.15in]{geometry}
\usepackage{amsmath,amssymb,amsthm}
\usepackage{mathtools}
\usepackage[T1]{fontenc}
\usepackage{hyperref}
\usepackage[numbers,sort&compress]{natbib}
\usepackage{url}
\usepackage{xcolor}
\usepackage{listings}
\usepackage{booktabs}
\usepackage{enumitem}
\usepackage{setspace}
\usepackage{fancyhdr}
\usepackage{titlesec}

\setlength{\headheight}{14pt}

% ------------------------------------------------------------------
% Theorem environments
% ------------------------------------------------------------------
\theoremstyle{plain}
\newtheorem{theorem}{Theorem}[section]
\newtheorem{lemma}[theorem]{Lemma}
\newtheorem{corollary}[theorem]{Corollary}
\newtheorem{proposition}[theorem]{Proposition}

\theoremstyle{definition}
\newtheorem{definition}[theorem]{Definition}
\newtheorem{conjecture}[theorem]{Conjecture}

\theoremstyle{remark}
\newtheorem{remark}[theorem]{Remark}

% ------------------------------------------------------------------
% Listings for Rocq
% ------------------------------------------------------------------
\definecolor{rocqblue}{rgb}{0.08,0.20,0.52}
\definecolor{rocqgray}{rgb}{0.45,0.45,0.45}
\definecolor{rocqgreen}{rgb}{0.08,0.42,0.08}
\definecolor{rocqbg}{rgb}{0.97,0.97,0.99}

\lstdefinestyle{rocq}{
  language=ML,
  backgroundcolor=\color{rocqbg},
  basicstyle=\ttfamily\footnotesize,
  keywordstyle=\color{rocqblue}\bfseries,
  commentstyle=\color{rocqgray}\itshape,
  stringstyle=\color{rocqgreen},
  numbers=left,
  numberstyle=\tiny\color{rocqgray},
  numbersep=8pt,
  frame=single,
  framerule=0.4pt,
  rulecolor=\color{rocqgray!40},
  breaklines=true,
  captionpos=b,
  morekeywords={Definition,Theorem,Lemma,Proof,Qed,Inductive,Record,
                Parameter,Axiom,Module,End,Type,forall,exists,fun,
                match,with,let,in,Prop,Set,Conjecture,Corollary,
                Section,Require,Import,Fixpoint,destruct,exact,
                apply,intros,unfold,right,left,split,simpl,
                reflexivity,auto,rewrite,classic,Forall},
  deletekeywords={not,and,or,true,false},
}

% ------------------------------------------------------------------
% Section title formatting
% ------------------------------------------------------------------
\titleformat{\section}{\large\bfseries}{\thesection.}{0.5em}{}
\titleformat{\subsection}{\normalsize\bfseries}{\thesubsection.}{0.5em}{}
\titleformat{\subsubsection}{\normalsize\itshape}{\thesubsubsection.}{0.5em}{}

% ------------------------------------------------------------------
% Header
% ------------------------------------------------------------------
\pagestyle{fancy}
\fancyhf{}
\rhead{\small\textit{Structural Admissibility in Rocq}}
\lhead{\small\textit{Anonymous}}
\cfoot{\thepage}
\renewcommand{\headrulewidth}{0.4pt}

% ------------------------------------------------------------------
% Macros
% ------------------------------------------------------------------
\newcommand{\simM}{\sim_M}
\newcommand{\simPhi}{\sim_\Phi}
\newcommand{\hatPhi}{\hat{\Phi}}
\newcommand{\WD}{\mathrm{WarrantDebt}}

% ==================================================================
\begin{document}
% ==================================================================

% ------------------------------------------------------------------
% Title block
% ------------------------------------------------------------------
\begin{titlepage}
\centering
\vspace*{2cm}

{\LARGE\bfseries Structural Admissibility in Rocq}\\[0.6em]
{\large Warrant Debt, Observational Quotients, and Type-Theoretic Verification}

\vspace{1.8cm}

{\large Anonymous}

\vspace{1.5cm}

{\normalsize Draft, \today}

\vspace{2.5cm}

\begin{abstract}
\noindent
Verification can fail before proof search begins. When the observation
map of a verification system conflates states that a correctness
predicate must distinguish, the verification question is malformed:
the required information has already been destroyed by the
observational quotient.

We formalise this condition as admissibility. A predicate is admissible
exactly when it factors through the equivalence relation induced by
observation, or equivalently when the observation-induced quotient is
fine enough to preserve the distinctions the predicate requires.

For perturbation systems, we introduce warrant debt as the operational
witness of admissibility failure: warrant debt occurs when algebraic
observation treats two perturbations as equivalent while empirical
evidence distinguishes them. The Warrant Debt Theorem is proved in
Rocq: whenever warrant debt exists, the corresponding empirical
tolerance predicate is not admissible.

The mechanisation shows that inadmissibility is a prior proof
obligation rather than a failure of proof search: verification must
establish that the correctness predicate survives quotienting before
proof search begins. The Rocq development fully proves the Warrant
Debt Theorem, the monitor-as-evidence and conditional safety results,
the concrete rupture decomposition theorem for the
\texttt{SimpleSynergyCompose} instance, and the Non-Factorisation
Theorem. The abstract composition theorem is stated at the interface
level and is provable only for a concrete composition instance
satisfying the required axioms. A higher-order extension to finite
dependency families is presented as a proposed generalisation and
is not part of the present mechanisation.
\end{abstract}
\vspace{0.8em}
\noindent\textbf{Keywords:} Formal verification, Observational equivalence,
Quotient semantics, Admissibility, Type theory, Rocq.
\end{titlepage}

\tableofcontents
\newpage

% ==================================================================
\section{Introduction}
\label{sec:intro}
% ==================================================================

The question is not whether a correctness predicate is true,
but whether truth survives quotienting by the observation structure.

Verification failure is commonly diagnosed as proof failure: the prover
did not find a certificate, the invariant was too weak, or the
abstraction was too coarse.
This diagnosis is often wrong.

Many failures are structural: the observation map has already collapsed
the distinctions the correctness predicate requires before any proof
attempt begins.
In that situation, proof search is not merely difficult.
The question it is attempting to answer is malformed.

When an observation map $M : X \to O$ is not injective, distinct states
share an observational image.
The induced equivalence
\[
x \sim_M y \iff M(x) = M(y)
\]
partitions the state space, and a predicate
\[
\Phi : X \to \{0,1\}
\]
is \emph{admissible} with respect to $M$ if and only if it factors
through that partition:
\[
\Phi = \hat{\Phi} \circ M
\]
for some
\[
\hat{\Phi} : O \to \{0,1\}.
\]

If $\Phi$ is not admissible, there exist states $x, y$ such that
$M(x) = M(y)$ and $\Phi(x) \neq \Phi(y)$.
No reasoning that operates within the observational regime can certify
$\Phi$, because the regime has erased the distinction $\Phi$ requires.
Inadmissibility is structural impossibility, not epistemic limitation.

The problem belongs to quotient semantics, abstract interpretation,
observational equivalence, and definability theory.
The issue is not the expressive weakness of the prover; it is whether
the observational quotient is fine enough to preserve the distinctions
the predicate demands.

Triple-overlap obstruction in regional algebra provides a concrete instance of inadmissibility: the non-vanishing Čech class is the witness that strict associativity is not representable within the existing gluing quotient.

In perturbation systems, this failure appears as
\emph{warrant debt}: algebraic observation treats two perturbations as
equivalent while empirical evidence distinguishes them.
Warrant debt is therefore the concrete witness that admissibility has
failed.

Dependent type theory makes admissibility constraints explicit as proof
obligations prior to theorem proving.
The question is not whether a proof can be found, but whether the
correctness predicate is even well-formed relative to the observation
structure.
The Rocq development presented here makes this precise: warrant debt
is represented as a formal object, and the theorem that its existence
implies non-admissibility is fully proved.

\paragraph{Contributions.}
Three results are established.
The Factorisation Theorem (Section~\ref{sec:quotients}) characterises
admissibility in three equivalent forms: factorisation through the
observation map, constancy on observation-induced equivalence classes,
and refinement of the correctness partition by the observational
partition.
It is the formal criterion against which all subsequent results are
measured.
The Warrant Debt Theorem (Section~\ref{sec:warrant}) proves that
whenever empirical evidence distinguishes algebraically equivalent
perturbations, the corresponding empirical tolerance predicate is not
admissible; warrant debt is thus the operational witness of
admissibility failure, and the theorem is fully mechanised in Rocq.
The mechanisation (Sections~\ref{sec:type}--\ref{sec:composition})
exhibits admissibility checking as a prior proof obligation preceding
monitor correctness and compositional reasoning. A concrete rupture
decomposition theorem is proved for the \texttt{SimpleSynergyCompose}
instance, supplying the binary compositional witness, and the
Non-Factorisation Theorem is established in the Rocq development. A
higher-order extension to finite dependency families on $n$
components is set out at the level of mathematical structure as a
proposed generalisation; minimal support sets identify the smallest
subsystems on which the global predicate genuinely depends. The
higher-order theorem is not part of the present mechanisation.

\paragraph{Scope.}
This is a formal verification paper.
Motivation is provided only where it identifies a formal problem.
Metaphysical interpretation belongs to a separate paper in this
programme.

% ==================================================================
\section{Observational Quotients and Admissibility}
\label{sec:quotients}
% ==================================================================

This section states the abstract theorem of which all subsequent
results are instances.
No perturbation-specific machinery is required here.

\subsection{Observation Maps and Observational Equivalence}

Let $X$ be a non-empty set (the state space) and $O$ a non-empty set
(the observation space).
Let
\[
M : X \to O
\]
be a function, the \emph{observation map}; $M$ need not be injective.

\begin{definition}[Observational Equivalence]
The \emph{observational equivalence} induced by $M$ is
\[
x \sim_M y
\;\stackrel{\mathrm{def}}{\iff}\;
M(x) = M(y).
\]
This is an equivalence relation on $X$.

Let $X/{\sim_M}$ denote the quotient and
\[
\pi : X \to X/{\sim_M}
\]
the canonical projection.

There exists a unique injection
\[
\bar{M} : X/{\sim_M} \to O
\]
such that
\[
M = \bar{M} \circ \pi.
\]
\end{definition}

The observational equivalence is the relevant object of study.
Two states that $M$ cannot separate are, from the perspective of the
verification system, identical.
No predicate that respects the observation structure can distinguish
them.

\subsection{Admissible Predicates}

\begin{definition}[Admissibility]
\label{def:admissibility}

Let
\[
\Phi : X \to \{0,1\}
\]
be a correctness predicate.

$\Phi$ is \emph{admissible with respect to $M$} if and only if there
exists
\[
\widetilde{\Phi} : X/{\sim_M} \to \{0,1\}
\]
such that
\[
\Phi = \widetilde{\Phi} \circ \pi.
\]

That is, $\Phi$ must factor through the quotient induced by
observational equivalence.

Equivalently, since
\[
M = \bar{M} \circ \pi,
\]
admissibility is equivalent to factorisation through the observation
map itself: there exists
\[
\hat{\Phi} : O \to \{0,1\}
\]
such that
\[
\Phi = \hat{\Phi} \circ M.
\]
\end{definition}

The quotient is primary.
The codomain $O$ matters only because the observation map induces the
equivalence relation that defines the quotient.
Verification is therefore governed not by the observation values
themselves, but by the partition they impose on the state space.

\subsection{The Factorisation Theorem}

\begin{theorem}[Factorisation]
\label{thm:factorisation}

The following conditions are equivalent.

\begin{enumerate}[label=(\arabic*),noitemsep]
  \item $\Phi$ is admissible with respect to $M$.

  \item $\Phi$ is constant on $\sim_M$-equivalence classes:
        \[
        M(x) = M(y)
        \;\Rightarrow\;
        \Phi(x) = \Phi(y).
        \]

  \item The observational equivalence refines the predicate
        equivalence:
        \[
        \sim_M \subseteq \sim_{\Phi},
        \]
        where
        \[
        x \sim_{\Phi} y
        \;\stackrel{\mathrm{def}}{\iff}\;
        \Phi(x) = \Phi(y).
        \]
\end{enumerate}
\end{theorem}

\begin{proof}

$(1)\Rightarrow(2)$:
If
\[
\Phi = \hat{\Phi} \circ M
\]
and
\[
M(x) = M(y),
\]
then
\[
\Phi(x)
=
\hat{\Phi}(M(x))
=
\hat{\Phi}(M(y))
=
\Phi(y).
\]

$(2)\Rightarrow(3)$:
If
\[
x \sim_M y,
\]
then
\[
M(x) = M(y),
\]
so by (2),
\[
\Phi(x) = \Phi(y),
\]
hence
\[
x \sim_{\Phi} y.
\]
Therefore
\[
\sim_M \subseteq \sim_{\Phi}.
\]

$(3)\Rightarrow(1)$:
Since
\[
\sim_M \subseteq \sim_{\Phi},
\]
the assignment
\[
\widetilde{\Phi}([x]_M) = \Phi(x)
\]
is well-defined on the quotient
\[
X/{\sim_M}.
\]
Thus
\[
\Phi = \widetilde{\Phi} \circ \pi.
\]

Extending $\widetilde{\Phi}$ along $\bar{M}$ yields the equivalent
factorisation through $O$.
\end{proof}

Theorem~\ref{thm:factorisation} is the spine of the paper.

A predicate is decidable from observation if and only if the
observation-induced partition is fine enough to respect the
predicate's partition.

If the observational equivalence is too coarse, no argument about the
predicate can recover what the quotient has already destroyed.

\subsection{Saturation}

\begin{definition}[Saturation]

A set
\[
S \subseteq X
\]
is $M$-\emph{saturated} if for all
\[
x \in S
\quad\text{and}\quad
y \in X,
\]
\[
M(x) = M(y)
\;\Rightarrow\;
y \in S.
\]

That is, $S$ is a union of $\sim_M$-equivalence classes.
\end{definition}

\begin{corollary}[Saturation Criterion]
\label{cor:saturation}

$\Phi$ is admissible with respect to $M$ if and only if
\[
\Phi^{-1}(1)
\]
is $M$-saturated.
\end{corollary}

\begin{proof}

Admissibility holds if and only if no $\sim_M$-equivalence class
straddles the boundary between
\[
\Phi^{-1}(1)
\quad\text{and}\quad
\Phi^{-1}(0).
\]

This holds if and only if
\[
\Phi^{-1}(1)
\]
is a union of $\sim_M$-classes, which is exactly the definition of
$M$-saturation.
\end{proof}

\subsection{Inadmissibility as Structural Impossibility}

When $\Phi$ is not admissible, there exist
\[
x, y \in X
\]
such that
\[
M(x) = M(y)
\quad\text{and}\quad
\Phi(x) \neq \Phi(y).
\]

The pair $(x,y)$ is an inadmissibility witness.

No predicate that factors through the current observation map can
recover this distinction while that observational regime remains fixed.
The quotient
\[
X/{\sim_M}
\]
has already collapsed the distinction that $\Phi$ requires.

Inadmissibility is not epistemic ignorance.
It is structural impossibility.

The failure lies in the representational regime itself, not in the
strength of the prover operating within it.

\subsection{Literature Placement}
\label{subsec:literature}

Theorem~\ref{thm:factorisation} is a definability result.
Its connections to the literature are structural.

In abstract interpretation~\cite{cousot1977,cousot1979}, a Galois
connection $(\alpha,\gamma)$ relates concrete and abstract domains;
a property survives abstraction only if the abstract domain is fine
enough to preserve the required distinctions~\cite{cousot2002,giacobazzi1998}.
The quotient $X/{\sim_M}$ is the abstract domain.
Inadmissibility is the case in which the abstraction is too coarse
and the required distinction has been irreversibly erased.

Beth's theorem~\cite{beth1953} establishes that implicit and explicit
definability coincide in first-order logic.
In the present setting there is a fixed observation map rather than a
first-order theory, and the observation-induced equivalence plays the
role of the theory's indistinguishability relation.
A predicate is definable in the relevant sense if and only if the
quotient is fine enough to preserve its partition~\cite{hodges1993}.

Two structures are elementarily equivalent if and only if the
Duplicator wins the Ehrenfeucht--Fra\"{i}ss\'{e} game of all finite
lengths~\cite{ehrenfeucht1961,fraisse1954}.
The relation $\sim_M$ is the observational analogue: two states are
$\sim_M$-equivalent when no observation separates them.
Inadmissibility is the corresponding structural result: collapse of
the observational quotient forces indistinguishability between states
that the correctness predicate must separate~\cite{libkin2004}.

Milner's observational equivalence~\cite{milner1980,milner1989}
identifies processes that no context can distinguish by observable
actions.
Observation maps play the role of contexts, and the admissibility
condition is the verification analogue of contextual
equivalence~\cite{sangiorgi2011}.

% ==================================================================
\section{Concrete Perturbation Space}
\label{sec:perturbations}
% ==================================================================

The abstract framework of Section~\ref{sec:quotients} is instantiated
here.

The observation map $M$ becomes an algebraic observation map over
perturbations: it preserves carrier class and magnitude while
discarding emergent boundary information.

The correctness predicates $\Phi$ become tolerance predicates over
interface regions.

This section introduces only the formal objects required for the
Warrant Debt Theorem of Section~\ref{sec:warrant}; nothing here is
independent of that purpose.

% ------------------------------------------------------------------
\subsection{Boundary Type Classification}
% ------------------------------------------------------------------

\begin{definition}[Boundary Types]

A perturbation is classified by one of three boundary types.

\begin{enumerate}[label=(\roman*),noitemsep]
    \item \texttt{INV} (invariant):
    tolerance conditions are fully determined at design time by
    algebraic structure;

    \item \texttt{REL} (relational):
    tolerance conditions vary within bounded and certifiable limits;

    \item \texttt{EMG} (emergent):
    tolerance conditions depend on information that cannot be
    recovered from the algebraic observation kernel alone.
\end{enumerate}

The distinction is representational rather than merely descriptive:
\texttt{EMG} marks precisely the case in which algebraic observation
is structurally insufficient.

\end{definition}

\begin{lstlisting}[style=rocq,
caption={Boundary types and perturbation record.}]
Inductive BType : Type :=
  | INV  (* invariant: structurally determined *)
  | REL  (* relational: bounded, variable      *)
  | EMG. (* emergent: observationally lost     *)

Record Perturbation := mkPerturbation {
  p_btype     : BType;
  p_carrier   : CarrierClass;
  p_magnitude : Interval;
  p_mode      : nat
}.
\end{lstlisting}

The \texttt{Interval} type carries a proof of well-formedness
$\ell \leq h$, so every magnitude is represented as a certified
interval rather than an uncertified real pair.

A structural consequence used later is that \texttt{EMG} is the top
element under composition: if either component is emergent, the
composition is emergent.
Composition cannot eliminate emergent character.

% ------------------------------------------------------------------
\subsection{Tolerance Predicates}
% ------------------------------------------------------------------

\begin{definition}[Tolerance Predicates]

A tolerance predicate specifies whether an interface region can absorb
a perturbation without violating its safety boundary.
Two forms are distinguished: formal tolerance, decided at design time
by a Boolean procedure; and empirical tolerance, indexed by evidence
such as sealed monitoring logs or certified test records.

\end{definition}

\begin{lstlisting}[style=rocq,
caption={Tolerance predicates.}]
Definition TolPredicate := Perturbation -> Prop.

Definition formal_tol (dec : Perturbation -> bool)
  : TolPredicate :=
  fun p => dec p = true.

Parameter Evidence : Type.
Parameter empirical_tol : Evidence -> TolPredicate.
\end{lstlisting}

The type \texttt{Evidence} is left abstract.
The results of Section~\ref{sec:warrant} do not depend on the internal
structure of evidence, only on the fact that empirical tolerance may
distinguish perturbations that algebraic observation treats as
equivalent.

This gap is the source of warrant debt.

% ------------------------------------------------------------------
\subsection{Rupture}
% ------------------------------------------------------------------

\begin{definition}[Rupture]
\label{def:rupture}

A perturbation $p$ \emph{ruptures} tolerance predicate
\texttt{tol} when

\[
\neg\,\mathtt{tol}(p).
\]

\end{definition}

\begin{lstlisting}[style=rocq,
caption={Rupture.}]
Definition Rupture
  (tol : TolPredicate)
  (p : Perturbation) : Prop :=
  ~ tol p.
\end{lstlisting}

Rupture is not a property of the perturbation alone.
It is a relation between a perturbation and a tolerance predicate.

Changing the predicate changes what counts as rupture.

This matters because detectability depends not only on whether rupture
occurs, but on whether the tolerance predicate itself is admissible
with respect to the observation structure.
Section~\ref{sec:warrant} makes this precise.

% ==================================================================
\section{Algebraic Observation and Warrant Debt}
\label{sec:warrant}
% ==================================================================

This section is the formal centre of the paper.

Section~\ref{sec:quotients} established the abstract admissibility
criterion: a predicate is admissible if and only if the observational
equivalence induced by the observation map is fine enough to preserve
the distinctions the predicate requires.

This section is the concrete instance of that criterion for
perturbation systems.

The algebraic observation map plays the role of $M$; \texttt{alg\_equiv}
plays the role of $\sim_M$; and \texttt{PredAdmissible} is the
perturbation-level factorisation condition of
Theorem~\ref{thm:factorisation}.

% ------------------------------------------------------------------
\subsection{The Algebraic Observation Map}
% ------------------------------------------------------------------

The algebraic observation map strips emergent boundary information.

Two perturbations are algebraically equivalent when they agree on
carrier class and magnitude and both belong to the algebraically
observable fragment of the system.

\begin{lstlisting}[style=rocq,
caption={Algebraic equivalence and predicate admissibility.}]
Definition btype_algebraic (b : BType) : bool :=
  match b with
  | INV => true
  | REL => true
  | EMG => false
  end.

Definition alg_equiv (p q : Perturbation) : Prop :=
  p_carrier p   = p_carrier q /\
  p_magnitude p = p_magnitude q /\
  btype_algebraic (p_btype p) = true /\
  btype_algebraic (p_btype q) = true.

Definition PredAdmissible
  (P : Perturbation -> Prop) : Prop :=
  forall p q : Perturbation,
    alg_equiv p q -> (P p <-> P q).
\end{lstlisting}

\begin{definition}[Observational Admissibility in Perturbation Space]
\label{def:pred-admissible}

A predicate $P$ is \emph{admissible} with respect to the algebraic
observation map if and only if it is invariant under
\texttt{alg\_equiv}: $P$ cannot distinguish perturbations that the
algebraic kernel cannot distinguish.

This is the perturbation-level realisation of condition~(2) of
Theorem~\ref{thm:factorisation}:
\[
M(x) = M(y)
\;\Rightarrow\;
\Phi(x) = \Phi(y).
\]

\end{definition}

Emergent perturbations lie outside the algebraic observation domain
entirely:
\[
\texttt{btype\_algebraic(p\_btype q) = false}
\]
whenever
\[
\texttt{p\_btype q = EMG}.
\]

The algebraic kernel therefore forms no equivalence class containing
an emergent perturbation.

Any predicate that separates an algebraic perturbation from an
emergent one witnesses information that the kernel destroys.

% ------------------------------------------------------------------
\subsection{Warrant Equivalence}
% ------------------------------------------------------------------

\begin{definition}[Warrant Equivalence]

Perturbations $p$ and $q$ are \emph{warrant-equivalent} with respect
to evidence $e$ when empirical tolerance treats them identically:

\[
\mathtt{warrant\_equiv}(e,p,q)
\;\stackrel{\mathrm{def}}{\iff}\;
\Bigl(
\mathtt{empirical\_tol}(e,p)
\leftrightarrow
\mathtt{empirical\_tol}(e,q)
\Bigr).
\]

\end{definition}

\begin{lstlisting}[style=rocq,
caption={Warrant equivalence.}]
Definition warrant_equiv
  (e : Evidence)
  (p q : Perturbation) : Prop :=
  empirical_tol e p <-> empirical_tol e q.
\end{lstlisting}

Warrant equivalence is the empirical counterpart of algebraic
equivalence.

Algebraic equivalence records what the formal certification
infrastructure treats as the same perturbation.

Warrant equivalence records what the empirical evidence treats as the
same perturbation.

Warrant debt arises precisely when these two notions of sameness
disagree.

% ------------------------------------------------------------------
\subsection{Warrant Debt}
% ------------------------------------------------------------------

\begin{definition}[Warrant Debt]
\label{def:warrant-debt}

A \emph{warrant debt} exists between perturbations $p$ and $q$ with
respect to evidence $e$ when they are algebraically equivalent but not
warrant-equivalent:

\[
\WD(e,p,q)
\;\stackrel{\mathrm{def}}{\iff}\;
\mathtt{alg\_equiv}(p,q)
\;\wedge\;
\neg\,\mathtt{warrant\_equiv}(e,p,q).
\]

\end{definition}

\begin{lstlisting}[style=rocq,
caption={Warrant debt.}]
Definition WarrantDebt
  (e : Evidence)
  (p q : Perturbation) : Prop :=
  alg_equiv p q /\ ~ warrant_equiv e p q.
\end{lstlisting}

A $\WD(e,p,q)$ is the formal witness that the observational regime is
too coarse.

The algebraic certification infrastructure treats $p$ and $q$ as
equivalent, while the evidence record distinguishes them.

Admissibility therefore fails at the level of representation itself:
the system is attempting to certify a predicate that its own
observation structure cannot support.

% ------------------------------------------------------------------
\subsection{The Warrant Debt Theorem}
\label{subsec:wdt}
% ------------------------------------------------------------------

\begin{theorem}[Warrant Debt Implies Non-Admissibility]
\label{thm:wdt}

For any evidence $e$ and perturbations $p$ and $q$,

\[
\WD(e,p,q)
\;\Rightarrow\;
\neg\,
\mathtt{PredAdmissible}
\bigl(
\mathtt{empirical\_tol}(e)
\bigr).
\]

\end{theorem}

\begin{proof}

Suppose
\[
\WD(e,p,q).
\]

Then
\[
\mathtt{alg\_equiv}(p,q)
\]
and
\[
\neg\,\mathtt{warrant\_equiv}(e,p,q).
\]

Assume for contradiction that

\[
\mathtt{PredAdmissible}
\bigl(
\mathtt{empirical\_tol}(e)
\bigr).
\]

By Definition~\ref{def:pred-admissible}, admissibility requires that
algebraically equivalent perturbations cannot be separated by the
predicate.

Hence
\[
\mathtt{alg\_equiv}(p,q)
\]
implies

\[
\mathtt{empirical\_tol}(e,p)
\leftrightarrow
\mathtt{empirical\_tol}(e,q),
\]

which is exactly

\[
\mathtt{warrant\_equiv}(e,p,q).
\]

This is the contradiction.
Admissibility forces warrant equivalence, while warrant debt asserts
its negation.
Both cannot hold simultaneously.

\end{proof}

The Rocq mechanisation is fully proved:

\begin{lstlisting}[style=rocq,
caption={Warrant Debt Theorem, fully proved in Rocq.}]
Theorem warrant_debt_implies_non_admissible :
  forall (e : Evidence)
         (p q : Perturbation),
  WarrantDebt e p q ->
  ~ PredAdmissible (empirical_tol e).
Proof.
  intros e p q [Hequiv Hne_warrant] Hadm.
  apply Hne_warrant.
  unfold PredAdmissible in Hadm.
  exact (Hadm p q Hequiv).
Qed.
\end{lstlisting}

The proof is four lines.
Its brevity reflects the precision of the definitions: the mathematics
is in the definitions, not in the proof scripts.
A correct formalisation makes the theorem transparent; a bloated one
obscures it.

The engineering consequence is immediate.
If a warrant debt can be exhibited, the system is operating with an
inadmissible observational regime.
More proof cannot repair representational failure.
The correct response is not stronger proof search, but refinement of
the observation map until admissibility is restored, or formal
recognition that the target predicate is not verifiable within the
current regime.

\paragraph{Connection to Theorem~\ref{thm:factorisation}.}

The Warrant Debt Theorem is the concrete instance of the Factorisation
Theorem for perturbation systems.

\texttt{alg\_equiv} is the induced equivalence $\sim_M$.

\texttt{PredAdmissible} is the factorisation condition.

The pair $(p,q)$ with $\WD(e,p,q)$ is the inadmissibility witness of
Theorem~\ref{thm:factorisation}: two states the observation map
conflates that the predicate must separate.

The implication of Theorem~\ref{thm:wdt} is the perturbation-level
form of condition~(2) in Theorem~\ref{thm:factorisation}.

These are not two theorems.
They are one theorem at two levels of abstraction.

% ==================================================================
\section{Type-Level Realisation}
\label{sec:type}
% ==================================================================

This section shows how admissibility appears as a proof obligation
inside dependent type theory.

The type structure does not replace the mathematics; it exhibits the
same obligations in machine-checkable form.
Refinement of observation corresponds to refinement of the formal
objects through which verification proceeds.

Theorem~\ref{thm:wdt} established that warrant debt implies
non-admissibility.
The present section states the operational consequence: before proving
a predicate, one must first establish that the predicate is
well-formed relative to the available observation structure.

% ------------------------------------------------------------------
\subsection{Observational Refinement as Product Extension}
% ------------------------------------------------------------------

When a predicate $\Phi$ is inadmissible with respect to an observation
map $M_i$, the corrective is not stronger proof search but
observational refinement: introduce an auxiliary observation that
separates the inadmissibility witnesses.

\begin{definition}[Observational Refinement]

Let
\[
M_i : X \to O_i.
\]

An \emph{observational refinement} of $M_i$ is a map

\[
M_{i+1} : X \to O_{i+1}
\]

such that

\[
x \sim_{M_{i+1}} y
\;\Rightarrow\;
x \sim_{M_i} y.
\]

That is, the refined observation induces a strictly finer partition.
Equivalently,

\[
M_{i+1} = (M_i, N_i)
\]

for some auxiliary observation

\[
N_i : X \to F_i,
\]

so that

\[
O_{i+1} \cong O_i \times F_i.
\]

\end{definition}

\begin{proposition}[Product Decomposition]
\label{prop:product}

If

\[
M_{i+1} = (M_i, N_i)
\]

is an observational refinement, then

\[
O_{i+1} \cong O_i \times F_i,
\]

and the projection

\[
\pi_1 : O_{i+1} \to O_i
\]

satisfies

\[
M_i = \pi_1 \circ M_{i+1}.
\]

The refined observation reduces to the original by projection.

\end{proposition}

Each refinement step therefore extends the type of the observation
space by a product factor.
The projection $\pi_1$ corresponds exactly to forgetting the new
factor.

Refinement is not replacement.
It is conservative extension: the old observation remains recoverable
by projection, while the new factor provides the separation required
to eliminate inadmissibility.

% ------------------------------------------------------------------
\subsection{Type Structure and Proof Obligations}
% ------------------------------------------------------------------

Type structure enforces the projection and refinement obligations
corresponding to Theorem~\ref{thm:factorisation}.

The issue is not whether a proof can be found, but whether the
predicate is even well-formed relative to the observation structure.

A predicate admissible with respect to $M_i$ can be stated without
appeal to distinctions outside the partition induced by $M_i$.
A non-admissible predicate cannot: any attempt to verify it requires
appealing to a finer partition, which is exactly the formal signal that
the observation map must be refined.

In the Rocq development, $\mathtt{PredAdmissible}(P)$ is the explicit
proof obligation that must be discharged before verification of a
tolerance predicate may proceed.
If it cannot be discharged, the observation map must be refined by
product extension as in Proposition~\ref{prop:product}, or it must be
formally certified that $P$ is not verifiable within the current
regime.
Only after admissibility is established does proof search for $P$
begin.

This ordering is not a recommendation about proof strategy.
It is a structural consequence of Theorem~\ref{thm:wdt}.

Warrant debt proves that proof search inside an inadmissible
observational regime is not incomplete but misdirected: the system is
attempting to certify a predicate its own representation cannot
support.

% ==================================================================
\section{Composition and Rupture Decomposition}
\label{sec:composition}
% ==================================================================

\subsection{Composition Signature}

Composition is expressed as a Rocq module type.
A compliant operator must satisfy three axioms: emergent type
propagation, magnitude monotonicity, and a bridge between composite
rupture and its sources.

\begin{lstlisting}[style=rocq,
caption={Abstract composition module signature.}]
Module Type COMPOSITION_SIG.
  Parameter par_compose :
    Perturbation -> Perturbation -> Perturbation.

  Parameter synergy_debt :
    Perturbation -> Perturbation -> Prop.

  Axiom compose_btype_emg :
    forall p q,
    p_btype p = EMG \/ p_btype q = EMG ->
    p_btype (par_compose p q) = EMG.

  Axiom compose_magnitude_ub :
    forall p q,
      (iv_hi (p_magnitude p) <= iv_hi (p_magnitude (par_compose p q)))%R /\
      (iv_hi (p_magnitude q) <= iv_hi (p_magnitude (par_compose p q)))%R.

  Axiom rupture_synergy_bridge :
    forall (tol : TolPredicate) (p q : Perturbation),
    Rupture tol (par_compose p q) ->
    Rupture tol p \/ Rupture tol q \/ synergy_debt p q.
End COMPOSITION_SIG.
\end{lstlisting}

\begin{definition}[Interaction Debt, Binary Case]
\label{def:synergy-debt}

Two perturbations incur \emph{interaction debt} (in this binary instance
also called \emph{synergy debt}) when both are individually tolerated
under a predicate, yet their composition ruptures under that same
predicate.
The binary case names the lowest non-trivial level of the interaction
debt hierarchy and is not reducible to component-level failure.
The general hierarchy, which classifies interaction debt by the size
of the minimal support set on which inadmissibility arises, is
developed in Section~\ref{subsec:higher-order}.

\end{definition}

\subsection{Proved Rupture Decomposition}

The concrete module \texttt{SimpleSynergyCompose} satisfies the axioms
of \texttt{COMPOSITION\_SIG} with full proofs.
The rupture decomposition is therefore established constructively.

\begin{theorem}[Rupture Decomposition]
\label{thm:rupture}

For any tolerance predicate $\mathtt{tol}$ and perturbations $p$, $q$,

\[
\mathrm{Rup}(\mathtt{tol},\, \mathtt{par\_compose}(p,q))
\;\Rightarrow\;
\mathrm{Rup}(\mathtt{tol}, p)
\;\vee\;
\mathrm{Rup}(\mathtt{tol}, q)
\;\vee\;
\mathtt{synergy\_debt}(p, q).
\]

\end{theorem}

\begin{lstlisting}[style=rocq,
caption={Proved rupture decomposition.}]
Lemma rupture_synergy_bridge :
  forall (tol : TolPredicate) (p q : Perturbation),
  Rupture tol (par_compose p q) ->
  Rupture tol p \/ Rupture tol q \/ synergy_debt p q.
Proof.
  intros tol p q Hrup.
  destruct (classic (tol p)) as [Hp | Hnp].
  - destruct (classic (tol q)) as [Hq | Hnq].
    + right. right. unfold synergy_debt. exists tol.
      exact (conj Hp (conj Hq Hrup)).
    + right. left. exact Hnq.
  - left. exact Hnp.
Qed.

Theorem rupture_decomposition_proved :
  forall (tol : TolPredicate) (p q : Perturbation),
  Rupture tol (SimpleSynergyCompose.par_compose p q) ->
  Rupture tol p \/ Rupture tol q \/
  (exists tol' : TolPredicate,
    tol' p /\ tol' q /\
    Rupture tol' (SimpleSynergyCompose.par_compose p q)).
Proof.
  intros tol p q Hrup.
  exact (SimpleSynergyCompose.rupture_synergy_bridge tol p q Hrup).
Qed.
\end{lstlisting}

The third disjunct identifies the interaction term: both components
satisfy a tolerance predicate, yet their composition violates it.
Interaction debt at the binary level is therefore not residual; it is
a named and typed obligation.

The role of Theorem~\ref{thm:rupture} is compositional.
Where Theorem~\ref{thm:wdt} identifies failure of admissibility at the
level of a single predicate, Theorem~\ref{thm:rupture} identifies the
sources of failure in composed systems for the binary case.
A composite rupture on two components decomposes into certified local
failures or a single binary interaction term.
The binary formulation is, however, only the first non-trivial case of
a hierarchy.
For systems with three or more components, a single residual
interaction term is no longer sufficient: it records that some debt
exists without identifying the minimal locus of failure.
Section~\ref{subsec:higher-order} develops the corresponding hierarchy.

% ------------------------------------------------------------------
\subsection{Granular Synergy and Higher-Order Warrant Debt}
\label{subsec:higher-order}
% ------------------------------------------------------------------

The binary formulation of rupture decomposition is only the first
non-trivial case.
In general, for a system with component family
$
C = \{C_1,\dots,C_n\}
$,
global inadmissibility cannot be reduced to a single residual
interaction term without loss of explanatory force.
Such a residual merely records that some interaction debt exists; it
does not identify the minimal locus of failure.

The correct generalisation is hierarchical.
For each non-empty subset
$
S \subseteq C
$,
let
$
\Sigma_S
$
denote the warrant debt associated with the interaction structure of
the components in $S$, namely the failure of the global predicate to
factor through the observation map restricted to that subsystem.
Then local failure is the case $|S| = 1$, pairwise interaction debt
is the case $|S| = 2$ recovered by Theorem~\ref{thm:rupture}, triple
interaction debt the case $|S| = 3$, and so forth.

Accordingly, global rupture admits the decomposition
\[
\Phi_{\mathrm{global}}
\;=\;
\bigvee_{S \subseteq C,\; S \neq \varnothing}
\Sigma_S,
\]
where the operative question is not merely whether some
$
\Sigma_S
$
is non-zero, but which subsets $S$ are minimal witnesses of
inadmissibility.

\begin{definition}[Dependency Family and Minimal Support]
\label{def:dependency-family}

The \emph{dependency family} of a global predicate $\Phi$ over the
component family $C$ is
\[
\mathcal{D}(\Phi)
\;=\;
\{\,S \subseteq C
\mid
\Phi \text{ requires } S \text{ for admissibility}\,\}.
\]
A subset $S \in \mathcal{D}(\Phi)$ is a \emph{minimal support set}
when no proper subset of $S$ lies in $\mathcal{D}(\Phi)$.
The minimal elements of $\mathcal{D}(\Phi)$ identify the smallest
subsystems on which $\Phi$ genuinely depends.

\end{definition}

A predicate may be admissible on every proper projection of size less
than $k$, yet fail to be admissible on the full system.
This is genuine $k$-th order synergy, not reducible to lower-order
interaction.

\begin{theorem}[Higher-Order Warrant Debt]
\label{thm:higher-order}

Let $\Phi$ be a global predicate on a component family
$C = \{C_1,\dots,C_n\}$ with admissibility-induced dependency family
$\mathcal{D}(\Phi)$.
If $\Phi$ is inadmissible relative to the observation map on $C$, and
if for every proper subset $S' \subsetneq S$ of size strictly less than
some $k \leq n$ the projected predicate $\Phi|_{S'}$ is admissible,
then there exists a minimal support set $S^* \in \mathcal{D}(\Phi)$
with $|S^*| \geq k$.
The set $S^*$ is the smallest subsystem on which higher-order warrant
debt arises, and it is the minimal witness of inadmissibility.

\end{theorem}

\begin{proof}[Proof sketch]
By the Factorisation Theorem (Section~\ref{sec:quotients}),
inadmissibility of $\Phi$ is equivalent to the existence of a pair of
global states identified by the observation map but separated by
$\Phi$.
Project that pair onto every subset $S \subseteq C$.
The collection of subsets on which the projected pair remains separated
forms an upward-closed family inside the powerset lattice of $C$;
its minimal elements are the minimal support sets of $\mathcal{D}(\Phi)$.
By hypothesis, no subset of size less than $k$ lies in this family, so
every minimal element has cardinality at least $k$.
Pick any such minimal element $S^*$.
By construction $\Phi|_{S^*}$ is inadmissible while $\Phi|_{S'}$ is
admissible for every $S' \subsetneq S^*$, so $S^*$ is a minimal
witness of inadmissibility at order $|S^*|$.
\end{proof}

\begin{remark}[Mechanisation status]
The binary case $k = 2$ is the content of Theorem~\ref{thm:rupture}
and is fully mechanised in Rocq via
\texttt{rupture\_synergy\_bridge} for the concrete
\texttt{SimpleSynergyCompose} instance.
The general case is established here at the level of mathematical
structure and is presented as a proposed extension to the present
development. A Rocq mechanisation parameterised over the cardinality
of $C$ would reduce to the lattice structure of the powerset of $C$,
but it is outside the scope of the present formalisation, which
proves only the binary concrete instance.
\end{remark}

\begin{remark}[Significance for industrial systems]
The significance is practical as well as formal.
In industrial systems, all local indicators may remain within
tolerance, and every pairwise relation may appear acceptable, while
the joint configuration violates a global invariant.
Pressure, temperature, compressor lag, and flow timing may each be
locally benign yet jointly inadmissible.
Such failures are invisible to anomaly detection understood as local
deviation, but they are precisely what admissibility theory detects.
\end{remark}

\subsection{Bridge Theorem and Audit Interpretation}
\label{subsec:bridge}

The bridge from rupture to audit therefore requires classification
rather than mere existential detection.
The bridge result must identify not only that interaction debt exists,
but the minimal support set on which it arises.
Only then can verification terminate in assignable warrant debt:
which observational boundary failed, which refinement is required,
and which institutional responsibility is implicated.

\begin{corollary}[Strengthened Rupture-to-Synergy Bridge]
\label{cor:bridge}

If a global predicate $\Phi$ on the component family $C$ ruptures, and
if no lower-order interaction debt of order strictly less than $k$
suffices to witness this rupture, then the rupture implies the
existence of a minimal support set $S^* \in \mathcal{D}(\Phi)$ of
order at least $k$ that witnesses higher-order warrant debt.

\end{corollary}

This is stronger than the binary bridge axiom of
Section~\ref{sec:composition}, which records existential interaction
debt without classifying its order.
The strengthened formulation converts a rhetorical remainder into a
formal object: the minimal witness of inadmissibility, located at a
definite level of the dependency hierarchy.

Together, the three results (Theorem~\ref{thm:wdt},
Theorem~\ref{thm:rupture}, and Theorem~\ref{thm:higher-order}) separate
distinct failure modes.
Of these, the first two are fully mechanised in Rocq for the concrete
\texttt{SimpleSynergyCompose} instance, while the third is established
here at the level of mathematical structure and is identified in
Section~\ref{sec:boundary} as a proposed extension of the Rocq
development.
Representational failure, certified by warrant debt, requires
refinement of the observation map.
Compositional failure, certified at the binary level by rupture
decomposition and at higher orders by minimal support sets within the
dependency family, is analysed within a fixed observational regime.
The distinction is structural: refinement cannot eliminate higher-order
warrant debt, and composition cannot repair inadmissibility.

% ==================================================================
\section{Mechanisation Boundary}
\label{sec:boundary}
% ==================================================================

The principal non-factorisation result is now fully proved in the Rocq
development.
In particular, the theorem \texttt{no\_alg\_factorisation},
established in \texttt{bdgi\_perturbation.v} and ported into the
canonical development file \texttt{bdgi\_perturbation\_proved.v},
establishes by explicit witness construction that the rupture
predicate cannot factor through the algebraic quotient induced by
\texttt{alg\_equiv}.
This discharges the earlier conjectural status of non-factorisation
and makes precise the claim that observational indistinguishability is
a structural rather than merely epistemic limitation.

% ------------------------------------------------------------------
\begin{theorem}[Non-Factorisation]
\label{thm:nonfact}

There is no algebraic homomorphism

\[
f :
\mathtt{Perturbation}
\to
\mathtt{AlgPerturbation}
\]

that preserves rupture for every tolerance predicate.

Equivalently, no total projection from the full perturbation space to
its algebraically observable fragment can preserve rupture behaviour
in general.

\end{theorem}

\begin{lstlisting}[style=rocq,
caption={Non-factorisation type signature.}]
Definition AlgPerturbation :=
  { p : Perturbation |
    btype_algebraic (p_btype p) = true }.

Definition AlgHom :=
  Perturbation -> AlgPerturbation.

Definition proj_alg (ap : AlgPerturbation) : Perturbation :=
  proj1_sig ap.

Definition RupturePreserving
  (f : AlgHom)
  (tol : TolPredicate) : Prop :=
  forall p : Perturbation,
    Rupture tol p <->
    Rupture tol (proj_alg (f p)).
\end{lstlisting}

The Rocq proof fixes an arbitrary $f : \mathtt{AlgHom}$ and an
emergent perturbation $p_0$ with
$\mathtt{p\_btype}(p_0) = \mathtt{EMG}$, then constructs a witness
tolerance predicate that holds on every algebraic perturbation and
fails at $p_0$.
This contradicts rupture preservation and discharges the proof
obligation explicitly within the canonical development file.

\begin{remark}[Correction of an earlier formulation]
\label{rem:emg-correction}
An earlier draft of this development presented a related claim,
labelled \texttt{emg\_sensitive\_not\_admissible}, asserting that any
predicate sensitive to the emergent--algebraic distinction is
inadmissible with respect to the algebraic observation map.
Under the present definition of \texttt{alg\_equiv} that formulation
is in fact false, and the corresponding lemma has been removed from
the development rather than retained as an \texttt{Admitted}
declaration.
The witness-based non-factorisation result stated above is the correct
theorem.
It captures the substantive content of the earlier claim, namely that
the algebraic regime cannot retain the distinctions on which rupture
depends, without the spurious universal quantification over
admissibility predicates that rendered the earlier formulation
unsound.
\end{remark}

% ------------------------------------------------------------------
\subsection{The Abstract Composition Boundary}
% ------------------------------------------------------------------

The only remaining admitted theorem in the development is the abstract
statement \texttt{rupture\_decomposition}, formulated over the
parameter \texttt{par\_compose}.
This is intentional.
Since \texttt{par\_compose} is introduced as an opaque composition
operator, the theorem is not derivable without additional composition
axioms.
Rather than proving an artefact of an unconstrained abstraction, the
development proves the concrete theorem
\texttt{rupture\_decomposition\_proved} for the explicit
\texttt{SimpleSynergyCompose} instance, where the required bridge laws
are established.

This distinction is methodological rather than provisional.
The abstract statement specifies the interface conditions required of
any admissible composition operator, and the concrete theorem
certifies one realised instance satisfying those conditions.
Any further module bound by \texttt{COMPOSITION\_SIG} obtains the
theorem at instantiation time, conditional on supplying the axioms the
signature names.

% ------------------------------------------------------------------
\subsection{Higher-Order Generalisation}
% ------------------------------------------------------------------

Higher-order warrant debt for systems with more than two interacting
components is not part of the present mechanisation.
The current Rocq development proves the binary case, in which
interaction debt first appears non-trivially.
The finite dependency lattice required for the general case is
mathematically straightforward but proof-theoretically lengthy, and is
reserved for subsequent work.

% ==================================================================
\section{Conclusion}
\label{sec:conclusion}
% ==================================================================

Formal verification does not certify truth first.
It certifies admissibility first.
Only after admissibility is established does proof begin.

The Factorisation Theorem states the abstract criterion: a predicate
is decidable from observation if and only if the observation-induced
partition is fine enough to preserve the distinctions the predicate
requires.

The Warrant Debt Theorem is the concrete operational instance of that
criterion.
When algebraically equivalent perturbations are distinguished by
empirical evidence, the corresponding empirical tolerance predicate is
not admissible.
This result is fully mechanised in Rocq.

The Rupture Decomposition Theorem extends the account to composed
systems.
At the binary level, composite failure reduces to one of two cases:
local component rupture or a named interaction term.
The binary result is fully proved in Rocq for the concrete
\texttt{SimpleSynergyCompose} instance, and the abstract theorem over
\texttt{COMPOSITION\_SIG} is provable for any module that supplies
the required composition axioms.

The higher-order extension generalises this account to the full
dependency family.
Global failures decompose into local failures, lower-order interaction
debts, or irreducible higher-order warrant debt determined by minimal
support sets within the dependency family of the global predicate.
The minimal support set identifies the smallest subsystem on which the
global predicate genuinely depends, and therefore the smallest locus
at which institutional responsibility can be assigned.
This higher-order extension is presented at the level of mathematical
structure and is identified as a proposed extension to the present
Rocq development rather than a result of it.

The Non-Factorisation Theorem is fully proved in the Rocq development
and rules out any total algebraic projection that preserves rupture
behaviour for every tolerance predicate.

Together, these results separate two distinct sources of verification
failure.
Representational failure occurs when the observational regime is too
coarse and admissibility fails; it requires refinement of the
observation map.
Compositional failure occurs when admissibility holds but interaction
produces rupture; it is analysed within the fixed observational regime.

More proof cannot repair the first.
More abstraction cannot eliminate the second.

The connections to abstract interpretation
\cite{cousot1977,cousot1979},
Ehrenfeucht--Fra\"{i}ss\'{e} indistinguishability
\cite{ehrenfeucht1961,libkin2004},
Beth definability
\cite{beth1953,hodges1993},
and observational equivalence
\cite{milner1980,milner1989}
are structural rather than decorative.

They place admissibility inside an established family of results
concerning abstraction, quotienting, and definability: when the
quotient destroys the distinction, proof cannot restore it.

One further observation.

Whitehead's account of process treats the constitution of a proposition
and the determination of its observational conditions as
inseparable~\cite{whitehead1929}.

The type
\[
\texttt{WarrantDebt}(e,p,q)
\]
gives this a precise formal referent.

It is the witness that the empirical record preserves a distinction
that the algebraic regime has excluded.

The evidence knows the difference.
The formal channel does not.

Verification is therefore not the reporting of pre-existing truth.
It is the construction, under explicit structural conditions, of the
proof obligations that make truth formally admissible in the first
place.

% ==================================================================
\section*{Statements and Declarations}
\addcontentsline{toc}{section}{Statements and Declarations}
% ==================================================================

\paragraph{Funding.}
The author received no funding, grants, or other financial support
for the research, authorship, or preparation of this manuscript.

\paragraph{Competing interests.}
The author has no competing financial or non-financial interests to
disclose that are relevant to the content of this article.

\paragraph{Data availability.}
This is a formal verification paper and contains no empirical
datasets. The Rocq source files corresponding to the mechanised
results reported in Sections~\ref{sec:warrant}--\ref{sec:boundary},
including the proofs of the Warrant Debt Theorem, the concrete
rupture decomposition theorem for \texttt{SimpleSynergyCompose}, and
the Non-Factorisation Theorem, are available in the canonical
development file \texttt{bdgi\_perturbation\_proved.v} at
\url{https://anonymous.4open.science/r/structural-admissibility-in-rocq-0C1C/}.

\paragraph{Code availability.}
The Rocq development is distributed under an open licence at the
anonymous repository cited above and compiles against Rocq 8.18 or later.

\paragraph{Author contributions.}
Omitted for blind review.

\paragraph{Ethics approval.}
Not applicable. The research did not involve human participants,
human tissue, or animals.

% ==================================================================
\bibliographystyle{plainnat}
\begin{thebibliography}{35}

\bibitem{abramsky1994}
S.~Abramsky and A.~Jung.
\newblock Domain theory.
\newblock In S.~Abramsky, D.~M. Gabbay, and T.~S.~E. Maibaum, editors,
  \emph{Handbook of Logic in Computer Science}, volume~3, pages 1--168.
  Clarendon Press, 1994.

\bibitem{barras1997}
B.~Barras et al.
\newblock The {Coq} proof assistant reference manual.
\newblock Technical report, INRIA, 1997. Version 6.1.

\bibitem{bertot2004}
Y.~Bertot and P.~Cast\'{e}ran.
\newblock \emph{Interactive Theorem Proving and Program Development:
  Coq'Art}.
\newblock Springer, 2004.

\bibitem{beth1953}
E.~W. Beth.
\newblock On {P}adoa's method in the theory of definition.
\newblock \emph{Indagationes Mathematicae}, 15:330--339, 1953.

\bibitem{clarke1981}
B.~L. Clarke.
\newblock A calculus of individuals based on `connection'.
\newblock \emph{Notre Dame Journal of Formal Logic}, 22(3):204--218, 1981.

\bibitem{cohn1997}
A.~G. Cohn, B.~Bennett, J.~Gooday, and N.~M. Gotts.
\newblock Qualitative spatial representation and reasoning with the region
  connection calculus.
\newblock \emph{GeoInformatica}, 1(3):275--316, 1997.

\bibitem{cousot1977}
P.~Cousot and R.~Cousot.
\newblock Abstract interpretation: A unified lattice model for static analysis
  of programs by construction or approximation of fixpoints.
\newblock In \emph{Proceedings of the 4th {ACM} Symposium on Principles of
  Programming Languages}, pages 238--252. {ACM}, 1977.

\bibitem{cousot1979}
P.~Cousot and R.~Cousot.
\newblock Systematic design of program analysis frameworks.
\newblock In \emph{Proceedings of the 6th {ACM} Symposium on Principles of
  Programming Languages}, pages 269--282. {ACM}, 1979.

\bibitem{cousot2002}
P.~Cousot.
\newblock Constructive design of a hierarchy of semantics of a transition
  system by abstract interpretation.
\newblock \emph{Theoretical Computer Science}, 277(1--2):47--103, 2002.

\bibitem{dijk1976}
E.~W. Dijkstra.
\newblock \emph{A Discipline of Programming}.
\newblock Prentice-Hall, 1976.

\bibitem{ehrenfeucht1961}
A.~Ehrenfeucht.
\newblock An application of games to the completeness problem for formalized
  theories.
\newblock \emph{Fundamenta Mathematicae}, 49:129--141, 1961.

\bibitem{fraisse1954}
R.~Fra\"{i}ss\'{e}.
\newblock Sur quelques classifications des syst\`{e}mes de relations.
\newblock \emph{Publications Scientifiques de l'Universit\'{e} d'Alger},
  1:35--182, 1954.

\bibitem{giacobazzi1998}
R.~Giacobazzi, F.~Ranzato, and F.~Scozzari.
\newblock Making abstract interpretations complete.
\newblock \emph{Journal of the {ACM}}, 47(2):361--416, 1998.

\bibitem{gonthier2008}
G.~Gonthier.
\newblock Formal proof --- the four-color theorem.
\newblock \emph{Notices of the {AMS}}, 55(11):1382--1393, 2008.

\bibitem{hales2015}
T.~Hales et al.
\newblock A formal proof of the {Kepler} conjecture.
\newblock \emph{Forum of Mathematics, Pi}, 5:e2, 2017.

\bibitem{hodges1993}
W.~Hodges.
\newblock \emph{Model Theory}.
\newblock Cambridge University Press, 1993.

\bibitem{jacobs1999}
B.~Jacobs.
\newblock \emph{Categorical Logic and Type Theory}.
\newblock Elsevier, 1999.

\bibitem{libkin2004}
L.~Libkin.
\newblock \emph{Elements of Finite Model Theory}.
\newblock Springer, 2004.

\bibitem{milner1980}
R.~Milner.
\newblock A calculus of communicating systems.
\newblock \emph{Lecture Notes in Computer Science}, 92, 1980.

\bibitem{milner1989}
R.~Milner.
\newblock \emph{Communication and Concurrency}.
\newblock Prentice Hall, 1989.

\bibitem{moore2024rsl}
Author.
\newblock Regional structure logic: Boundary closure, primitive connection, and
  structural admissibility.
\newblock Manuscript, 2024.

\bibitem{nipkow2002}
T.~Nipkow, L.~C. Paulson, and M.~Wenzel.
\newblock \emph{Isabelle/{HOL}: A Proof Assistant for Higher-Order Logic}.
\newblock Springer, 2002.

\bibitem{pierce2002}
B.~C. Pierce.
\newblock \emph{Types and Programming Languages}.
\newblock MIT Press, 2002.

\bibitem{randell1992}
D.~A. Randell, Z.~Cui, and A.~G. Cohn.
\newblock A spatial logic based on regions and connection.
\newblock In \emph{Proceedings of the 3rd International Conference on
  Knowledge Representation and Reasoning}, pages 165--176.
  Morgan Kaufmann, 1992.

\bibitem{sangiorgi2011}
D.~Sangiorgi.
\newblock \emph{Introduction to Bisimulation and Coinduction}.
\newblock Cambridge University Press, 2011.

\bibitem{simons1987}
P.~Simons.
\newblock \emph{Parts: A Study in Ontology}.
\newblock Clarendon Press, 1987.

\bibitem{smith2000}
B.~Smith and A.~C. Varzi.
\newblock Fiat and bona fide boundaries.
\newblock \emph{Philosophy and Phenomenological Research},
  60(2):401--420, 2000.

\bibitem{spivak2014}
D.~I. Spivak.
\newblock \emph{Category Theory for the Sciences}.
\newblock MIT Press, 2014.

\bibitem{stell2000}
J.~G. Stell.
\newblock Boolean connection algebras: A new approach to the region connection
  calculus.
\newblock \emph{Artificial Intelligence}, 122(1--2):111--136, 2000.

\bibitem{tabareau2021}
N.~Tabareau, E.~Tanter, and M.~Sozeau.
\newblock The marriage of univalence and parametricity.
\newblock \emph{Journal of the {ACM}}, 68(1):1--44, 2021.

\bibitem{whitehead1929}
A.~N. Whitehead.
\newblock \emph{Process and Reality: An Essay in Cosmology}.
\newblock Macmillan, 1929.

\bibitem{winskel1993}
G.~Winskel.
\newblock \emph{The Formal Semantics of Programming Languages}.
\newblock MIT Press, 1993.

\end{thebibliography}

\end{document}
